% !Mode:: "TeX:UTF-8"
% !TEX root = ../../main.tex
% ============================================================
% TOPIC 18
% ============================================================
\mashgulot{18}{Operations on linear operators. The space of linear operators. The matrix of a linear operator under a change of basis}

\begin{tarif}
A mapping $\mathcal A\colon V\to V$ of a linear space $V$ into itself is
called a \textbf{linear operator} if for all $x,y\in V$ and all $\lambda\in\R$
\[\kf{\mathcal A(x+y)=\mathcal Ax+\mathcal Ay,\qquad \mathcal A(\lambda x)=\lambda\,\mathcal Ax .}\]
\end{tarif}

\begin{tarif}
Let $e_1,\dots,e_n$ be a basis and $\mathcal Ae_j=\sum_{i=1}^na_{ij}e_i$. The matrix
$A=(a_{ij})$ is called the \textbf{matrix of the operator} in this basis:
its $j$th column consists of the coordinates of $\mathcal Ae_j$. The
coordinates of $y=\mathcal Ax$ are then found from the formula $Y=AX$.
\end{tarif}

\begin{misol}
Write down the matrix of the operator $\mathcal A(x_1;x_2)=(2x_1-x_2;\ x_1+3x_2)$
in the standard basis and find the image of $x=(1;-2)$.

\textbf{Solution.} Since $\mathcal Ae_1=\mathcal A(1;0)=(2;1)$ and
$\mathcal Ae_2=\mathcal A(0;1)=(-1;3)$, we get
$A=\begin{pmatrix}2&-1\\1&3\end{pmatrix}$ and
$\mathcal Ax=\begin{pmatrix}2&-1\\1&3\end{pmatrix}\begin{pmatrix}1\\-2\end{pmatrix}=\begin{pmatrix}4\\-5\end{pmatrix}$.
\end{misol}

\subsection{Operations on linear operators}

Let $A$ and $B$ be the matrices of operators $\mathcal A$ and $\mathcal B$ in a
fixed basis. Then $\mathcal A+\mathcal B$ has the matrix $A+B$, and $\lambda\mathcal A$
has the matrix $\lambda A$. The product $\mathcal A\mathcal B$ (first $\mathcal B$
is applied, then $\mathcal A$) has the matrix $AB$. Finally, the inverse operator
$\mathcal A^{-1}$ exists if and only if $\det A\neq0$, and its matrix is $A^{-1}$.
Consequently, the set of linear operators on an $n$-dimensional space is itself a
linear space; it is isomorphic to the space of $n\times n$ matrices and has dimension $n^2$.

\begin{misol}
Let $\mathcal R$ be the rotation of the plane through $90^\circ$, with matrix
$R=\begin{pmatrix}0&-1\\1&0\end{pmatrix}$, and let $\mathcal S$ be the reflection in
the $Ox$ axis, $S=\begin{pmatrix}1&0\\0&-1\end{pmatrix}$. Then
\[\kf{SR=\begin{pmatrix}0&-1\\-1&0\end{pmatrix},\qquad
RS=\begin{pmatrix}0&1\\1&0\end{pmatrix}.}\]
Since $SR\neq RS$, multiplication of operators is not commutative.
\end{misol}

\subsection{The matrix of an operator in a new basis}

\begin{teorema}
Let $T$ be the transition matrix from a basis $e$ to a basis $f$ (the columns of
$T$ are the coordinates of the vectors $f_j$ in the basis $e$). Then the matrix of
the operator $\mathcal A$ in the basis $f$ is
\[\kf{A'=T^{-1}AT .}\]
Similar matrices have the same determinant, trace and characteristic polynomial.
\end{teorema}

\begin{misol}
In the basis $e$ an operator has the matrix $A=\begin{pmatrix}1&2\\0&3\end{pmatrix}$.
Find its matrix $A'$ in the new basis $f_1=e_1+e_2$, $f_2=e_1$.

\textbf{Solution.} Here $T=\begin{pmatrix}1&1\\1&0\end{pmatrix}$ and
$T^{-1}=\begin{pmatrix}0&1\\1&-1\end{pmatrix}$, so
\[\kf{A'=T^{-1}AT=\begin{pmatrix}0&1\\1&-1\end{pmatrix}
\begin{pmatrix}1&2\\0&3\end{pmatrix}\begin{pmatrix}1&1\\1&0\end{pmatrix}
=\begin{pmatrix}3&0\\0&1\end{pmatrix}.}\]
In the new basis the matrix of the operator is diagonal: $\mathcal Af_1=3f_1$ and
$\mathcal Af_2=f_2$. As a check, $\det A=\det A'=3$ and $\operatorname{tr}A=\operatorname{tr}A'=4$.
\end{misol}

\begin{tarif}
The set $\ker\mathcal A=\{x\colon\mathcal Ax=\theta\}$ is called the \textbf{kernel}
of the operator, and $\operatorname{Im}\mathcal A=\{\mathcal Ax\colon x\in V\}$ its
\textbf{image}; they satisfy $\dim\ker\mathcal A+\dim\operatorname{Im}\mathcal A=\dim V$.
\end{tarif}

\subsection{Applications of linear operators in information security}

\textbf{The encryption operator.} The Hill cipher is the linear operator
$\mathcal K(p)=Kp\pmod{26}$ on the space $\mathbb Z_{26}^n$. Decryption is possible
only if this operator is invertible, i.e.\ $\gcd(\det K,26)=1$. Two successive
encryptions $\mathcal K_2\mathcal K_1$ again form a Hill cipher, with key $K_2K_1$;
for this reason “double encryption” does not make the cipher any stronger.

\textbf{The linear layers of AES.} The \emph{ShiftRows} and \emph{MixColumns}
steps of AES are linear operators over $\mathrm{GF}(2^8)$. They provide
“diffusion”, while the nonlinear \emph{SubBytes} step provides “confusion”. A
cipher built from linear layers alone is easily broken by linear algebra.

\begin{misol}
Encrypting first with $K_1=\begin{pmatrix}1&1\\0&1\end{pmatrix}$ and then with
$K_2=\begin{pmatrix}2&0\\1&1\end{pmatrix}$ is the same as encrypting once with the key
$K=K_2K_1=\begin{pmatrix}2&2\\1&2\end{pmatrix}$.
\end{misol}

\subsection{Computing with Python, Maple and Mathcad}

The change of basis in the example above is a single matrix product. In SymPy:
\begin{Verbatim}[fontsize=\footnotesize,frame=single,label={\small Python}]
from sympy import Matrix
A = Matrix([[1, 2], [0, 3]])
T = Matrix([[1, 1], [1, 0]])
B = T.inv() * A * T
print(B)
print(A.det(), B.det(), A.trace(), B.trace())
\end{Verbatim}
\begin{Verbatim}[fontsize=\footnotesize,frame=single,label={\small Output}]
Matrix([[3, 0], [0, 1]])
3 3 4 4
\end{Verbatim}
In Maple the matrix product is written with a dot:
\begin{Verbatim}[fontsize=\footnotesize,frame=single,label={\small Maple}]
with(LinearAlgebra):
A := Matrix([[1, 2], [0, 3]]):
T := Matrix([[1, 1], [1, 0]]):
B := MatrixInverse(T) . A . T;
Determinant(B);  Trace(B);
\end{Verbatim}
In Mathcad the inverse is entered as the power $-1$, and \texttt{tr(B)} gives the trace:
\begin{Verbatim}[fontsize=\footnotesize,frame=single,label={\small Mathcad}]
A := | 1  2 |      T := | 1  1 |
     | 0  3 |           | 1  0 |
B := T^-1.A.T        B = | 3  0 |
                          | 0  1 |
|B| = 3      tr(B) = 4
\end{Verbatim}

\subsection{Exercises on linear operators}

\emph{1. The matrices $A$ and $B$ of two operators are given. Find the matrices of the indicated operators.}
\bigskip
\begin{enumerate}[label=\arabic*., leftmargin=*]
\item $A=\begin{pmatrix}-2 & -3\\ 2 & 0\end{pmatrix}$, $B=\begin{pmatrix}-1 & -2\\ -2 & 0\end{pmatrix}$; \ $A+3B$, $AB$, $BA$
\item $A=\begin{pmatrix}0 & -2\\ 0 & -1\end{pmatrix}$, $B=\begin{pmatrix}0 & -1\\ -2 & -1\end{pmatrix}$; \ $A+3B$, $AB$, $BA$
\item $A=\begin{pmatrix}-3 & -1\\ 1 & 2\end{pmatrix}$, $B=\begin{pmatrix}3 & 3\\ -2 & 3\end{pmatrix}$; \ $A-B$, $AB$, $BA$
\item $A=\begin{pmatrix}-2 & -2\\ -2 & 2\end{pmatrix}$, $B=\begin{pmatrix}-1 & 1\\ 1 & -2\end{pmatrix}$; \ $A+3B$, $AB$, $BA$
\item $A=\begin{pmatrix}3 & 3\\ 3 & 0\end{pmatrix}$, $B=\begin{pmatrix}-2 & 2\\ -1 & -1\end{pmatrix}$; \ $A-2B$, $AB$, $BA$
\item $A=\begin{pmatrix}-1 & 0\\ -2 & 3\end{pmatrix}$, $B=\begin{pmatrix}2 & 0\\ 1 & -1\end{pmatrix}$; \ $A+2B$, $AB$, $BA$
\item $A=\begin{pmatrix}1 & 1\\ -1 & 3\end{pmatrix}$, $B=\begin{pmatrix}-1 & 1\\ 2 & -1\end{pmatrix}$; \ $A+2B$, $AB$, $BA$
\item $A=\begin{pmatrix}-3 & -1\\ 0 & 1\end{pmatrix}$, $B=\begin{pmatrix}1 & -2\\ 0 & 2\end{pmatrix}$; \ $A+3B$, $AB$, $BA$
\item $A=\begin{pmatrix}3 & -2\\ -2 & -2\end{pmatrix}$, $B=\begin{pmatrix}2 & 2\\ 3 & 0\end{pmatrix}$; \ $A-B$, $AB$, $BA$
\item $A=\begin{pmatrix}0 & 1\\ 0 & 1\end{pmatrix}$, $B=\begin{pmatrix}-3 & 2\\ -2 & 1\end{pmatrix}$; \ $A+2B$, $AB$, $BA$
\item $A=\begin{pmatrix}3 & -1\\ 3 & 3\end{pmatrix}$, $B=\begin{pmatrix}3 & 3\\ 2 & 1\end{pmatrix}$; \ $A+B$, $AB$, $BA$
\item $A=\begin{pmatrix}-3 & 0\\ -2 & 0\end{pmatrix}$, $B=\begin{pmatrix}-2 & 0\\ -3 & -2\end{pmatrix}$; \ $A+3B$, $AB$, $BA$
\item $A=\begin{pmatrix}-3 & 1\\ -1 & -2\end{pmatrix}$, $B=\begin{pmatrix}3 & 2\\ 0 & 3\end{pmatrix}$; \ $A+2B$, $AB$, $BA$
\item $A=\begin{pmatrix}3 & 1\\ 2 & 3\end{pmatrix}$, $B=\begin{pmatrix}1 & -3\\ -2 & 3\end{pmatrix}$; \ $A+3B$, $AB$, $BA$
\item $A=\begin{pmatrix}3 & -3\\ 0 & 0\end{pmatrix}$, $B=\begin{pmatrix}3 & 0\\ -2 & 0\end{pmatrix}$; \ $A+B$, $AB$, $BA$
\item $A=\begin{pmatrix}-2 & -1 & 2\\ 0 & 2 & 2\\ -1 & -2 & -1\end{pmatrix}$, $B=\begin{pmatrix}-1 & 2 & 3\\ 1 & -1 & -2\\ 0 & 2 & 1\end{pmatrix}$; \ $A-2B$, $AB$, $BA$
\item $A=\begin{pmatrix}1 & 0 & -2\\ 3 & -2 & 1\\ 2 & 3 & 2\end{pmatrix}$, $B=\begin{pmatrix}0 & -3 & 1\\ 1 & 1 & 3\\ 0 & 1 & -3\end{pmatrix}$; \ $A-2B$, $AB$, $BA$
\item $A=\begin{pmatrix}3 & 1 & 3\\ -1 & 0 & 2\\ 0 & 0 & 1\end{pmatrix}$, $B=\begin{pmatrix}1 & 0 & 3\\ 0 & 3 & 3\\ 3 & -2 & -2\end{pmatrix}$; \ $A-B$, $AB$, $BA$
\item $A=\begin{pmatrix}3 & 3 & -2\\ -1 & 2 & 2\\ -1 & 0 & -3\end{pmatrix}$, $B=\begin{pmatrix}3 & 3 & -3\\ -1 & 3 & 1\\ 2 & -2 & 3\end{pmatrix}$; \ $A+2B$, $AB$, $BA$
\item $A=\begin{pmatrix}-2 & -2 & -1\\ -3 & -2 & 3\\ 1 & 2 & 3\end{pmatrix}$, $B=\begin{pmatrix}-2 & -2 & -1\\ -3 & 1 & -1\\ -3 & -1 & -1\end{pmatrix}$; \ $A-2B$, $AB$, $BA$
\item $A=\begin{pmatrix}1 & -1 & -1\\ 2 & -2 & 0\\ 2 & 0 & 3\end{pmatrix}$, $B=\begin{pmatrix}1 & 0 & 2\\ 1 & 2 & -3\\ -1 & 0 & 3\end{pmatrix}$; \ $A+3B$, $AB$, $BA$
\item $A=\begin{pmatrix}-3 & 1 & -1\\ 0 & -1 & -3\\ -2 & -1 & -2\end{pmatrix}$, $B=\begin{pmatrix}-1 & -2 & 2\\ 2 & 2 & 0\\ -3 & -3 & 3\end{pmatrix}$; \ $A+3B$, $AB$, $BA$
\item $A=\begin{pmatrix}-1 & 3 & 3\\ 1 & 3 & 0\\ 1 & 1 & -3\end{pmatrix}$, $B=\begin{pmatrix}-2 & 2 & -3\\ 3 & 0 & 3\\ -1 & -2 & 3\end{pmatrix}$; \ $A+3B$, $AB$, $BA$
\item $A=\begin{pmatrix}-3 & -2 & -1\\ 3 & 1 & 2\\ 0 & 1 & 3\end{pmatrix}$, $B=\begin{pmatrix}1 & 3 & 0\\ 3 & 0 & 2\\ -1 & 0 & 3\end{pmatrix}$; \ $A-2B$, $AB$, $BA$
\item $A=\begin{pmatrix}1 & 2 & 3\\ -2 & 2 & -2\\ -3 & -3 & 0\end{pmatrix}$, $B=\begin{pmatrix}2 & 3 & -1\\ 2 & -3 & 1\\ -3 & -1 & -2\end{pmatrix}$; \ $A-B$, $AB$, $BA$
\item $A=\begin{pmatrix}1 & -2 & 2\\ 2 & -1 & -3\\ -2 & -2 & 2\end{pmatrix}$, $B=\begin{pmatrix}2 & 0 & -1\\ -3 & 1 & 0\\ 2 & 2 & 1\end{pmatrix}$; \ $A-B$, $AB$, $BA$
\item $A=\begin{pmatrix}-1 & 0 & -2\\ 0 & 1 & -2\\ 1 & -1 & 0\end{pmatrix}$, $B=\begin{pmatrix}-3 & 1 & 1\\ 2 & 3 & -1\\ 1 & -3 & 1\end{pmatrix}$; \ $A+B$, $AB$, $BA$
\item $A=\begin{pmatrix}-1 & 3 & 2\\ -1 & 2 & 1\\ -3 & -2 & 2\end{pmatrix}$, $B=\begin{pmatrix}2 & 3 & 1\\ 1 & -2 & -3\\ -2 & -1 & 3\end{pmatrix}$; \ $A-B$, $AB$, $BA$
\item $A=\begin{pmatrix}-2 & 0 & -3\\ 3 & 1 & 3\\ -2 & -3 & 0\end{pmatrix}$, $B=\begin{pmatrix}1 & 2 & -2\\ 0 & -3 & 2\\ 2 & 3 & 0\end{pmatrix}$; \ $A+B$, $AB$, $BA$
\item $A=\begin{pmatrix}-2 & -2 & 3\\ 0 & -1 & 2\\ -1 & 0 & -1\end{pmatrix}$, $B=\begin{pmatrix}-2 & -2 & 0\\ 3 & 2 & -1\\ 2 & 2 & 2\end{pmatrix}$; \ $A+2B$, $AB$, $BA$
\end{enumerate}

\emph{2. The matrix $A$ of an operator in a basis $e$ and the transition matrix $T$ are given. Find the matrix $A^{\prime}=T^{-1}AT$ in the new basis.}
\bigskip
\begin{multicols}{2}\raggedright
\begin{enumerate}[label=\arabic*., leftmargin=*]
\item $A=\begin{pmatrix}4 & 0\\ 1 & -1\end{pmatrix}$, $T=\begin{pmatrix}0 & -1\\ 1 & -2\end{pmatrix}$
\item $A=\begin{pmatrix}-3 & 1\\ -2 & 1\end{pmatrix}$, $T=\begin{pmatrix}-1 & -2\\ 0 & -1\end{pmatrix}$
\item $A=\begin{pmatrix}-1 & -2\\ -2 & 1\end{pmatrix}$, $T=\begin{pmatrix}1 & 0\\ 1 & 1\end{pmatrix}$
\item $A=\begin{pmatrix}1 & 4\\ -3 & 0\end{pmatrix}$, $T=\begin{pmatrix}1 & 0\\ -1 & 1\end{pmatrix}$
\item $A=\begin{pmatrix}0 & -1\\ 3 & -2\end{pmatrix}$, $T=\begin{pmatrix}1 & 0\\ 0 & 1\end{pmatrix}$
\item $A=\begin{pmatrix}1 & 1\\ 4 & 1\end{pmatrix}$, $T=\begin{pmatrix}-1 & 1\\ -1 & 0\end{pmatrix}$
\item $A=\begin{pmatrix}4 & 3\\ -3 & 3\end{pmatrix}$, $T=\begin{pmatrix}1 & 1\\ -1 & 0\end{pmatrix}$
\item $A=\begin{pmatrix}-2 & 2\\ 2 & -4\end{pmatrix}$, $T=\begin{pmatrix}2 & -1\\ 3 & -2\end{pmatrix}$
\item $A=\begin{pmatrix}-2 & 4\\ 1 & 1\end{pmatrix}$, $T=\begin{pmatrix}1 & 0\\ 1 & -1\end{pmatrix}$
\item $A=\begin{pmatrix}2 & 4\\ 2 & -2\end{pmatrix}$, $T=\begin{pmatrix}1 & -2\\ 0 & -1\end{pmatrix}$
\item $A=\begin{pmatrix}4 & -1\\ 3 & 3\end{pmatrix}$, $T=\begin{pmatrix}0 & 1\\ -1 & -2\end{pmatrix}$
\item $A=\begin{pmatrix}0 & -4\\ -4 & -3\end{pmatrix}$, $T=\begin{pmatrix}-1 & 0\\ -2 & 1\end{pmatrix}$
\item $A=\begin{pmatrix}0 & 2\\ -2 & 3\end{pmatrix}$, $T=\begin{pmatrix}-1 & 3\\ 0 & -1\end{pmatrix}$
\item $A=\begin{pmatrix}-2 & -3\\ 0 & 2\end{pmatrix}$, $T=\begin{pmatrix}2 & 1\\ 1 & 1\end{pmatrix}$
\item $A=\begin{pmatrix}-2 & 3\\ 4 & 3\end{pmatrix}$, $T=\begin{pmatrix}1 & -1\\ 1 & -2\end{pmatrix}$
\item $A=\begin{pmatrix}-2 & -2\\ 0 & 2\end{pmatrix}$, $T=\begin{pmatrix}1 & 0\\ 1 & -1\end{pmatrix}$
\item $A=\begin{pmatrix}4 & -2\\ 2 & 2\end{pmatrix}$, $T=\begin{pmatrix}0 & 1\\ -1 & 0\end{pmatrix}$
\item $A=\begin{pmatrix}1 & 0\\ 1 & 0\end{pmatrix}$, $T=\begin{pmatrix}-1 & 0\\ 1 & 1\end{pmatrix}$
\item $A=\begin{pmatrix}-1 & -4\\ -4 & -3\end{pmatrix}$, $T=\begin{pmatrix}2 & -1\\ 3 & -2\end{pmatrix}$
\item $A=\begin{pmatrix}4 & 4\\ -1 & 3\end{pmatrix}$, $T=\begin{pmatrix}0 & 1\\ 1 & -1\end{pmatrix}$
\item $A=\begin{pmatrix}0 & -2\\ 4 & -1\end{pmatrix}$, $T=\begin{pmatrix}0 & 1\\ -1 & -1\end{pmatrix}$
\item $A=\begin{pmatrix}-3 & -3\\ -2 & -4\end{pmatrix}$, $T=\begin{pmatrix}3 & -2\\ 2 & -1\end{pmatrix}$
\item $A=\begin{pmatrix}-2 & 0\\ 0 & 2\end{pmatrix}$, $T=\begin{pmatrix}0 & 1\\ -1 & 3\end{pmatrix}$
\item $A=\begin{pmatrix}1 & 3\\ 3 & 4\end{pmatrix}$, $T=\begin{pmatrix}-1 & -2\\ -1 & -1\end{pmatrix}$
\item $A=\begin{pmatrix}0 & 2\\ -4 & -1\end{pmatrix}$, $T=\begin{pmatrix}-1 & 0\\ -2 & 1\end{pmatrix}$
\item $A=\begin{pmatrix}-4 & -1\\ -2 & 0\end{pmatrix}$, $T=\begin{pmatrix}-1 & 0\\ 1 & 1\end{pmatrix}$
\item $A=\begin{pmatrix}3 & 1\\ 0 & -2\end{pmatrix}$, $T=\begin{pmatrix}1 & 1\\ 0 & -1\end{pmatrix}$
\item $A=\begin{pmatrix}-4 & -3\\ -3 & -4\end{pmatrix}$, $T=\begin{pmatrix}1 & 0\\ 0 & 1\end{pmatrix}$
\item $A=\begin{pmatrix}-4 & 0\\ -4 & -1\end{pmatrix}$, $T=\begin{pmatrix}-2 & -1\\ 3 & 2\end{pmatrix}$
\item $A=\begin{pmatrix}3 & 3\\ 2 & 1\end{pmatrix}$, $T=\begin{pmatrix}1 & -1\\ -1 & 2\end{pmatrix}$
\end{enumerate}
\end{multicols}

\emph{3. Applied problems (10 problems).}
\bigskip
\begin{enumerate}[label=\arabic*., leftmargin=*]
\item \textbf{(Linearity.)} Is the operator $\mathcal A(x_1;x_2)=(x_1+2x_2;\,3x_1-x_2)$ linear? Write down its matrix.
\item \textbf{(A nonlinear map.)} Is the operator $\mathcal A(x_1;x_2)=(x_1+1;\,x_2)$ linear?
\item \textbf{(Rotation.)} Write down the matrix of the rotation of the plane through $90^\circ$ and find the image of $\vec x=(3;1)$.
\item \textbf{(Projection.)} Write down the matrix of the projection of $\mathbb R^3$ onto the $Oxy$ plane and verify that $\mathcal A^2=\mathcal A$.
\item \textbf{(Reflection.)} Write down the matrix of the reflection in the $Ox$ axis and verify that $\mathcal A^2=E$.
\item \textbf{(Image of a vector.)} Find the image of $\vec x=(1;-1;2)$ under the operator with matrix $A=\left(\begin{smallmatrix}1&2&0\\0&1&-1\\2&0&3\end{smallmatrix}\right)$.
\item \textbf{(Kernel.)} Find the kernel and the image of the operator with matrix $A=\left(\begin{smallmatrix}1&2\\2&4\end{smallmatrix}\right)$.
\item \textbf{(Composition.)} Find the matrix of the operator that first rotates the plane through $90^\circ$ and then reflects it in the $Ox$ axis.
\item \textbf{(Inverse operator.)} Find the matrix of the inverse of the operator with matrix $A=\left(\begin{smallmatrix}2&1\\5&3\end{smallmatrix}\right)$.
\item \textbf{(The Hill cipher as an operator.)} Let $\mathcal K(p)=Kp\pmod{26}$ with $K=\left(\begin{smallmatrix}3&2\\5&7\end{smallmatrix}\right)$. Show that $\mathcal K$ is invertible on $\mathbb Z_{26}^2$.
\end{enumerate}

